\section{Related Work}
\label{sec:related}

\textbf{LLM-as-a-judge.}  Zheng et al.~\cite{zheng2023judging} popularised
strong LLMs as judges of open-ended answers (MT-Bench, Chatbot Arena) and
documented position, verbosity and self-enhancement biases.  Wang et
al.~\cite{wang2023fair} showed that the ranking of two responses can be
manipulated simply by changing their order of appearance and proposed
calibration strategies.  Gu et al.~\cite{gu2024survey} survey the design and
reliability of LLM judges.  We adopt the pairwise protocol and the
order-swapping analysis of this line of work.

\textbf{Judging code.}  ICE-Score~\cite{zhuo2024icescore} instructs LLMs to
assess code usefulness and functional correctness without test oracles, and
CodeJudge~\cite{tong2024codejudge} guides LLMs to ``slow thinking'' to judge
semantic correctness of generated code without test cases.
\CJB{}~\cite{jiang2025codejudgebench} is the largest benchmark of LLM judges for
coding: 5{,}352 pairs over code generation, code repair and unit-test generation
on LiveCodeBench-v6~\cite{jain2024livecodebench}, with 26 judges.  It finds that
thinking models are much stronger judges, that pairwise comparison beats
point-wise scoring (point-wise scores tie about half of the time), that raw
responses with comments and reasoning judge better than extracted code, and that
judges are sensitive to response order.  All of its instances are single-file
and test-decidable; \RJB{} moves the same protocol to repository-level diffs
whose defects tests do not catch.

\textbf{Repository-level benchmarks and contamination.}
SWE-bench~\cite{jimenez2024swebench} contains 2{,}294 issue--pull-request
instances from 12 Python repositories.  SWE-Bench+~\cite{aleithan2024swebenchplus}
audited it and found widespread solution leakage and weak tests.
SWE-bench-Live~\cite{zhang2025swebenchlive} continuously collects new
issue-resolution tasks (initially 1{,}319 tasks from issues created since 2024
across 93 repositories), each with a Docker image built by its RepoLaunch
pipeline; we build on its released tasks and images.  LiveCodeBench~\cite{jain2024livecodebench} similarly avoids contamination by
collecting problems over time windows and evaluating on problems released after
a model's training cutoff.

\textbf{Execution-free verifiers for software agents.}  R2E-Gym~\cite{jain2025r2egym}
compares execution-based and execution-free verifiers and reports that the
latter are biased and often rely on stylistic features.
SWE-RM~\cite{shum2025swerm} trains execution-free reward models for software
engineering agents, and R4P~\cite{xu2025r4p} verifies patches by group-wise
reasoning, reporting 72.2\% accuracy on SWE-bench patches.  Li~\cite{li2026frozen}
shows that fixed LLM judges make version-dependent errors when comparing
coding-agent upgrades.  Agent-as-a-Judge~\cite{zhuge2024agent} instead lets
agentic systems evaluate agentic systems.

\textbf{Code review.}  CodeReviewer~\cite{li2022codereviewer} pre-trains a model
for code-change quality estimation, review generation and code refinement on
review data from nine programming languages; quality estimation is the closest
pre-LLM formulation of our task.  SWE-PRBench~\cite{kumar2026sweprbench}
evaluates AI code review on 350 pull requests against human-flagged issues and
reports that models detect only 15--31\% of them and that adding more context
degrades results.

\textbf{Mutation testing.}  Our variants are mutants in the sense of mutation
testing~\cite{jia2011mutation}; unlike classical mutants, they are
\emph{selected} to survive the test suite and to instantiate a
review-detectable defect class.
